\documentclass[./main.tex]{subfiles}
\begin{document}

% Slide # 3
\begin{frame}[t, label=slide03]
        % Title
        \frametitle{Dynamical systems that tip}

        % Body
        \footnotesize

        Consider $n-$dimensional Itô processes with additive diffusion $\sigma = \sqrt{2D}>0$ and non-autonomous drift

        \begin{equation*}
           \begin{cases}
                   dX_{t} = f(X_{t},\,\mu(t))dt + \sqrt{2D}dW_{t} \,, \\
              \dot{\mu} \in \Lambda_{\varepsilon}(\mu_{-},\mu_{+}) \,. 
           \end{cases}
        \end{equation*}

        \vspace{0.5cm}
        \uncover<2->{
                The law of evolution of the parameter is a\alr{parametrised family of shifts}in time connecting the past-limit (i.e. $\mu(t)\to\mu_{-}$, $t\to-\infty$) to the future-limit system (i.e. $\mu(t)\to\mu_{+}$, $t\to+\infty$).
        }

        \vspace{0.5cm}
        \uncover<3->{
                For any given $\mu$ an\alr{equilibrium}of the system is a state $X^{*}\in \mathbb{R}^{n}$ s.t. $f(X^{*},\mu) = 0$. 
        }

        \vspace{0.5cm}
        \uncover<4->{
                Since $\mu = \mu(t)$ then the equilbrium itself drifts in state space, i.e. $X^{*} = X_{t}^{*}$ is a\alr{quasi-steady equilibrium (QSE)}.
        }

        \vspace{0.5cm}
        \uncover<5->{
                If at any point $\mu(t)$ a QSE is created, destroyed or it changes its stability then a\alr{bifurcation}has occured.
        }
\end{frame}

\end{document}
